% Documentclass Settings
	\documentclass[a4paper, twoside]{article}
	% Margins
	\usepackage{geometry}
		\geometry{
			% showframe,
			left   = 20mm,
			right  = 20mm,
			top    = 20mm,
			bottom = 20mm
		}

	% Headers and footers
	\usepackage{fancyhdr}
		\pagestyle{fancy}
		\fancyhf{}
		\fancyhead[LE, RO]{\thepage}
		\fancyhead[CE]{\textsc{Cong Wen}}
		\fancyhead[CO]{\textsc\rightmark}

		\setlength{\headheight}{15pt}
		\renewcommand\headrulewidth{0pt}








\input{../universal-pre}

\title{\tbf{Park2022}}
\author{\tbf{Notes} by Cong Wen, 02/2023}
\date{}
% Last Updated: July 11, 2022



\begin{document}
\maketitle
\thispagestyle{empty}
\begin{abstract}
	This is my self-use notes\footnote{Last updated on February 07, 2023} on the talk given by \tbf{Sung Gi Park} at \tbf{MIT} on his work\cite{Park2022BaseChange}. Please reach out if you find anything incorrect.

\end{abstract}
\tableofcontents
% ----------------------------------------------------------------------------------------------------

One of the main results of \cite{Park2022BaseChange} is the following, generalizing a classical theorem of Viehweg-Zuo\cite[Theorem~0.2]{VZ2001} in case $n=1$.
\begin{thm}[{\cite[Corollary~1.10]{Park2022BaseChange}}]
	Let $X$ be a complex smooth projective variety of $\kpa(X)\geq0$, $f:X\ar\bP^{n}$ a surjective morphism, $\Dta(f)\sbs X$ the discriminant locus, then
	\[
		\Dim\Dta(f)=n-1,\Deg\Dta(f)\geq n+2.
	\]
\end{thm}


One of the key points of the paper is an improvement of Viehweg's fiber product trick in the logarithm setting to study the positivity of the pushforward of pluricanonical bundles.










% ----------------------------------------------------------------------------------------------------
\bibliographystyle{alpha}
\bibliography{../universal-ref}
\end{document}

